% !TeX root = ../paper.tex
\subsection{Results Interpretation and Analysis}
\label{sec:interpretation}

% TODO: Outline framework for interpreting evaluation results

\subsubsection{Metric-based Interpretation}
\label{sec:metric-interpretation}

\paragraph{Diagnostic Performance Interpretation}
% TODO: Define interpretation criteria for accuracy, F1, AUROC
% TODO: Analyze precision vs. recall trade-offs

\paragraph{Comparative Performance Analysis}
Interpret performance differences across methods:

\begin{itemize}
    \item \textbf{RAG vs. Fine-tuning:}
    \begin{itemize}
        \item If RAG outperforms: External knowledge is crucial; dataset may be limited
        \item If Fine-tuning outperforms: Sufficient training data; learned patterns are key
        \item If similar: Both knowledge and learned patterns contribute equally
    \end{itemize}
    
    \item \textbf{Hybrid vs. Individual Methods:}
    \begin{itemize}
        \item Positive fusion gain: Methods are complementary
        \item No gain: Methods capture similar information (redundancy)
        \item Negative impact: Integration introduces noise or confusion
    \end{itemize}
    
    \item \textbf{Performance by Disease Category:}
    \begin{itemize}
        \item Better on common diseases: Model relies on training frequency
        \item Better on rare diseases: External knowledge (RAG) is effective
        \item Consistent across categories: Robust generalization
    \end{itemize}
\end{itemize}


\subsubsection{Method-specific Analysis}
\label{sec:method-analysis}

\paragraph{RAG System Analysis}
Interpret RAG-specific results:

\begin{itemize}
    \item \textbf{Retrieval Quality Impact:}
    \begin{itemize}
        \item High Precision@K but low final accuracy: Generation weakness, not retrieval
        \item Low Precision@K but high accuracy: Model compensates for poor retrieval
        \item Both high: Optimal RAG pipeline
    \end{itemize}
    
    \item \textbf{Faithfulness vs. Accuracy:}
    \begin{itemize}
        \item High faithfulness, high accuracy: Retrieved knowledge is relevant and used correctly
        \item High faithfulness, low accuracy: Retrieved knowledge is insufficient or incorrect
        \item Low faithfulness: Model ignores retrieved context (may indicate poor integration)
    \end{itemize}
    
    \item \textbf{Knowledge Source Analysis:}
    \begin{itemize}
        \item Which sources (textbooks, articles, guidelines) contribute most to correct diagnoses?
        \item Are certain disease categories better supported by specific sources?
    \end{itemize}
\end{itemize}

\paragraph{Fine-tuned Model Analysis}
Interpret fine-tuning results:

\begin{itemize}
    \item \textbf{Training Dynamics:}
    \begin{itemize}
        \item Converged training: Model learned effectively
        \item Overfitting (train $\gg$ val): Need regularization or more data
        \item Underfitting (both high): Insufficient model capacity or training
    \end{itemize}
    
    \item \textbf{LoRA Impact:}
    \begin{itemize}
        \item Compare full fine-tuning vs. LoRA (if both tested)
        \item Assess parameter efficiency vs. performance trade-off
    \end{itemize}
    
    \item \textbf{Domain Adaptation Success:}
    \begin{itemize}
        \item Improvement over base model indicates successful adaptation
        \item Analyze which aspects (vocabulary, reasoning, multimodal fusion) improved most
    \end{itemize}
\end{itemize}

\paragraph{Hybrid System Analysis}
Interpret combination results:

\begin{itemize}
    \item \textbf{Complementarity:}
    \begin{itemize}
        \item Cases where RAG succeeds and fine-tuning fails (and vice versa)
        \item Identify patterns indicating when each method is preferred
    \end{itemize}
    
    \item \textbf{Fusion Strategy Effectiveness:}
    \begin{itemize}
        \item Compare different fusion methods (ensemble, gating, sequential)
        \item Assess computational cost vs. performance gain
    \end{itemize}
\end{itemize}


\subsubsection{Dataset and Data Quality Analysis}
\label{sec:dataset-analysis}

Interpret results in the context of data characteristics:

\begin{itemize}
    \item \textbf{Data Sufficiency:}
    \begin{itemize}
        \item Learning curves: Does performance plateau or still improving?
        \item Sample size analysis: Which disease categories need more data?
    \end{itemize}
    
    \item \textbf{Data Quality Impact:}
    \begin{itemize}
        \item Correlation between data quality (completeness, accuracy) and performance
        \item Impact of noisy labels or missing modalities
    \end{itemize}
    
    \item \textbf{Modality Contribution:}
    \begin{itemize}
        \item Ablation results: Which modalities (text, images, structured data) contribute most?
        \item Are certain diseases better diagnosed with specific modalities?
    \end{itemize}
    
    \item \textbf{Distribution Analysis:}
    \begin{itemize}
        \item Class imbalance effects: Performance on rare vs. common diseases
        \item Temporal/geographic distribution: Does model generalize across contexts?
    \end{itemize}
\end{itemize}


\subsubsection{Error Pattern Analysis}
\label{sec:error-patterns}

Systematic reasoning about model failures:

\begin{itemize}
    \item \textbf{Common Error Types:}
    \begin{itemize}
        \item Confusion between similar diseases (e.g., malaria vs. dengue)
        \item Missing rare diseases with subtle presentations
        \item Over-reliance on common symptoms leading to frequent diagnoses
    \end{itemize}
    
    \item \textbf{Root Causes:}
    \begin{itemize}
        \item \textbf{Data limitations:} Insufficient examples, imbalanced distribution, poor quality
        \item \textbf{Model limitations:} Inadequate reasoning, multimodal fusion issues, knowledge gaps
        \item \textbf{Task complexity:} Inherently ambiguous cases, multiple valid diagnoses
    \end{itemize}
    
    \item \textbf{Failure Mode Categorization:}
    \begin{itemize}
        \item Knowledge failures: Missing critical medical knowledge
        \item Reasoning failures: Incorrect logical inference
        \item Perception failures: Misinterpretation of images or text
        \item Integration failures: Poor multimodal fusion
    \end{itemize}
\end{itemize}


\subsubsection{Clinical Implications}
\label{sec:clinical-implications}

Interpret results from a clinical perspective:

\begin{itemize}
    \item \textbf{Diagnostic Safety:}
    \begin{itemize}
        \item Critical error analysis: Could misdiagnoses cause patient harm?
        \item Uncertainty calibration: Does model appropriately express confidence?
        \item Safe deployment: Which clinical scenarios are appropriate for system use?
    \end{itemize}
    
    \item \textbf{Clinical Utility:}
    \begin{itemize}
        \item Decision support value: How much does system improve clinician performance?
        \item Workflow integration: Can system be practically deployed in clinical settings?
        \item Time efficiency: Does system reduce time to diagnosis?
    \end{itemize}
    
    \item \textbf{Reliability and Trust:}
    \begin{itemize}
        \item Consistency: Stable performance across diverse cases?
        \item Explainability: Are diagnostic reasoning explanations clinically sound?
        \item Limitations awareness: Does system recognize its boundaries?
    \end{itemize}
\end{itemize}


\subsubsection{Comparative Reasoning Framework}
\label{sec:comparative-reasoning}

Structured reasoning about comparative results:

\begin{enumerate}
    \item \textbf{Observation:} State the quantitative findings
    \begin{itemize}
        \item Example: "RAG achieved F1=0.82, Fine-tuning F1=0.78, Hybrid F1=0.87"
    \end{itemize}
    
    \item \textbf{Interpretation:} Explain what the numbers mean
    \begin{itemize}
        \item Example: "Hybrid approach shows 5-9\% improvement, indicating complementary strengths"
    \end{itemize}
    
    \item \textbf{Reasoning:} Why did this happen?
    \begin{itemize}
        \item Example: "RAG excels on rare diseases via external knowledge, while fine-tuned model captures common diagnostic patterns. Hybrid leverages both."
    \end{itemize}
    
    \item \textbf{Evidence:} Support with qualitative analysis
    \begin{itemize}
        \item Example: "Error analysis shows RAG fails on diseases absent from knowledge base, while fine-tuning struggles with rare conditions in training set"
    \end{itemize}
    
    \item \textbf{Implications:} What does this mean for the research?
    \begin{itemize}
        \item Example: "Suggests combination architectures are essential for comprehensive diagnostic coverage"
    \end{itemize}
\end{enumerate}


\subsubsection{Hypothesis Validation}
\label{sec:hypothesis-validation}

Interpret results in relation to initial hypotheses (Section~\ref{sec:problem-statement}):

\begin{itemize}
    \item \textbf{H1:} RAG outperforms standalone fine-tuning for rare diseases
    \begin{itemize}
        \item Supported/Rejected: [Based on your results]
        \item Reasoning: [Explain why based on metrics and analysis]
    \end{itemize}
    
    \item \textbf{H2:} Hybrid approach achieves superior performance
    \begin{itemize}
        \item Supported/Rejected: [Based on your results]
        \item Reasoning: [Explain why based on metrics and analysis]
    \end{itemize}
    
    \item \textbf{H3:} Multimodal training improves generalization
    \begin{itemize}
        \item Supported/Rejected: [Based on your results]
        \item Reasoning: [Explain why based on metrics and analysis]
    \end{itemize}
\end{itemize}


\subsubsection{Limitations and Future Directions}
\label{sec:limitations-interpretation}

Critical analysis of approach limitations revealed through evaluation:

\begin{itemize}
    \item \textbf{Data Limitations:}
    \begin{itemize}
        \item Dataset size, diversity, quality constraints
        \item Impact on model generalization and robustness
    \end{itemize}
    
    \item \textbf{Model Limitations:}
    \begin{itemize}
        \item Architectural constraints
        \item Computational resource limitations
        \item Scalability concerns
    \end{itemize}
    
    \item \textbf{Evaluation Limitations:}
    \begin{itemize}
        \item Test set representativeness
        \item Metric comprehensiveness
        \item Lack of prospective clinical validation
    \end{itemize}
    
    \item \textbf{Future Improvements:}
    \begin{itemize}
        \item Enhanced architectures or training strategies indicated by analysis
        \item Additional data sources or modalities
        \item Improved evaluation frameworks
        \item Clinical deployment pathways
    \end{itemize}
\end{itemize}
